\label{ch:background}

\section{Autoregressive language-model serving}

An autoregressive language model generates one token at a time. Serving a
request normally has two computational phases. During \emph{prefill}, the
model processes the prompt and constructs the key--value cache. During
\emph{decode}, the service repeatedly executes the model to produce one or a
small number of new tokens. Prefill has substantial parallel work over prompt
tokens, whereas decode is iterative and often sensitive to memory access,
scheduling delay, and interference.

Modern serving engines batch requests whose sequences are at different
positions. vLLM introduced PagedAttention to reduce fragmentation and enable
flexible sharing of KV-cache blocks, thereby increasing the number of requests
that can be served concurrently \cite{kwon2023vllm}. Continuous batching makes
resource utilization higher, but also means that a background operation may
affect multiple requests in different phases at the same time.

Four metrics are central in this thesis. Output throughput is the number of
generated tokens per second over the measured interval. \gls{ttft} is the time
from request submission until its first output token, so it includes admission,
queueing, and prefill. \gls{tpot} averages the gaps between output tokens after
the first token and is primarily associated with decode. End-to-end latency
measures the complete request. Reporting both median and P99 values separates
typical behavior from tail behavior. Precise definitions are given in
\cref{app:metrics}.

\section{Sparse mixture-of-experts models}

The sparsely gated MoE layer was developed as a practical form of conditional
computation \cite{shazeer2017moe}. Let $x$ denote a token representation,
$E_i$ the $i$th expert, and $g_i(x)$ a router weight. A simplified MoE output is

\begin{equation}
  y(x)=\sum_{i\in\operatorname{TopK}(g(x))} g_i(x) E_i(x).
\end{equation}

Only the selected experts execute for $x$. This decouples total parameter
capacity from active arithmetic. GShard demonstrated large-scale conditional
computation with automatic sharding \cite{lepikhin2020gshard}; Switch
Transformer simplified routing to one expert per token and highlighted both
the scaling opportunity and the communication and load-balancing challenges
\cite{fedus2021switch}. Recent open models continue this direction. Qwen3
includes dense and MoE variants \cite{yang2025qwen3}, while DeepSeek-V3 uses a
large sparse architecture and an auxiliary-loss-free training-time balancing
strategy \cite{deepseek2024v3}.

The evaluated Qwen3-30B-A3B-Instruct-2507 checkpoint has 30.5 billion total
parameters, approximately 3.3 billion activated parameters, 48 layers, 128
experts, and eight selected experts per token \cite{qwen3modelcard}. The gap
between total and active parameters makes it a useful serving model: compute
is sparse, but expert weights remain substantial objects that must be placed
and occasionally copied.

\section{Expert parallelism}

In \gls{ep}, experts are distributed across ranks. A token is dispatched to
the ranks holding its selected experts; expert outputs are then returned to
the rank that continues the model computation. EP avoids replicating every
expert on every device, but introduces all-to-all-style communication and a
dependence on balanced rank workloads.

Let $a_{l,e}$ be the number of token assignments to logical expert $e$ in layer
$l$ over an observation window. If rank $r$ hosts a set $H_{l,r}$, a simple
load estimate is

\begin{equation}
  L_{l,r}=\sum_{e\in H_{l,r}} a_{l,e}.
\end{equation}

Even if training encourages globally balanced routing, a finite serving window
can be skewed. Request language, task, domain, and generation state can all
change which experts are selected. Moreover, the system cares about load per
physical rank, not only load per logical expert. Several moderately popular
experts placed together can overload a rank.

\section{Expert Parallel Load Balancing in vLLM}

vLLM's EPLB records expert usage over a configurable window and periodically
computes a new placement. Redundant experts add physical slots beyond the
number required for one copy of each logical expert. The balancer can use those
slots to replicate hot experts and distribute their assignments. vLLM exposes
the observation window, rebalance interval, redundancy, asynchronous mode,
policy, and migration communicator through \code{EPLBConfig}
\cite{vllmeplbdocs}.

Placement is represented by replicated logical-to-physical and
physical-to-logical mappings for every MoE layer. Once a new map is available,
each rank can independently derive which logical experts it loses, keeps, or
gains. Local rearrangements copy weights between slots on the same rank.
Remote gains require a source rank that already owns the logical expert and a
point-to-point weight transfer.

Synchronous and asynchronous EPLB have different performance objectives. A
synchronous rearrangement pauses inference. In that setting, extra sequential
communication rounds directly extend the pause, so quickly completing the
migration is the natural goal. Asynchronous EPLB moves experts in a background
worker while inference continues. It transfers one layer and coordinates its
activation with forward progress before proceeding. Here, the migration has
more time, but interference with foreground work matters. This thesis changes
only the asynchronous path and retains the synchronous behavior.

\begin{figure}[t]
  \centering
  \begin{tikzpicture}[
    node distance=13mm and 11mm,
    box/.style={draw,rounded corners,align=center,minimum width=31mm,minimum height=10mm},
    arr/.style={-{Latex[length=2mm]},thick}
  ]
    \node[box] (metrics) {Per-layer expert\\usage window};
    \node[box,right=of metrics] (placement) {EPLB policy\\new placement};
    \node[box,right=of placement] (schedule) {Migration\\schedule};
    \node[box,below=of schedule] (transfer) {Background layer\\transfer};
    \node[box,left=of transfer] (activate) {Mapping swap at\\safe boundary};
    \node[box,left=of activate] (serve) {Foreground\\inference};
    \draw[arr] (metrics) -- (placement);
    \draw[arr] (placement) -- (schedule);
    \draw[arr] (schedule) -- (transfer);
    \draw[arr] (transfer) -- (activate);
    \draw[arr] (activate) -- (serve);
    \draw[arr] (serve) -- (metrics);
    \draw[arr,dashed]
      (serve.south) to[out=-65,in=-115]
      node[below]{shared network resources} (transfer.south);
  \end{tikzpicture}
  \caption{Simplified asynchronous EPLB control loop. The contribution of this
  thesis is the migration-scheduling block and its execution path.}
  \label{fig:eplb-loop}
\end{figure}

\section{NIXL communication}

\gls{nixl} provides a transfer abstraction for distributed inference across
heterogeneous memory and storage. A transfer agent manages registered memory,
remote metadata, backend selection, and asynchronous point-to-point requests
\cite{nixlguide}. vLLM supports NIXL as an EPLB communicator and can also use
Gloo or other backends depending on mode and configuration
\cite{vllmconfigdocs}.

The final evaluation explicitly selects NIXL. This choice matters for both
performance and correctness. In the implementation evaluated here,
\code{execute()} drains the queued transfer operations and clears the transfer
context. Consequently, every migration batch must restore the layer context,
post only that batch's sends and receives, and call \code{execute()}. Moving
the call outside the loop would collapse the batch boundaries and defeat the
scheduler's concurrency constraint.
